\chapter{Kết luận và hướng phát triển}
\section{Kết quả đạt được}
EduShield AI đã tích hợp quy trình thi trắc nghiệm từ tài khoản giảng viên, nhập đề, công bố lịch, phê duyệt sinh viên đến làm bài, chấm điểm và xuất kết quả. Phiên thi được cấp theo lượt ghi danh, giúp ràng buộc quyền tham gia và xác định danh tính dùng trong nghiệp vụ sau phê duyệt.

Mốc bắt đầu lưu tại máy chủ và điều kiện hạn nộp bảo vệ thời gian làm bài; ràng buộc duy nhất ngăn bản nộp trùng. Vi phạm được gắn với đề và lượt ghi danh, hỗ trợ truy vết và thống kê theo đúng bài thi. Hệ thống có dọn ảnh, sao lưu, kiểm tra toàn vẹn và công cụ phục hồi dữ liệu.

Phần AI tích hợp Face Landmarker và EfficientDet Lite0 ngay trong trình duyệt. Các quy tắc hình học, độ sáng, ưu tiên nhãn và cooldown biến đầu ra mô hình thành cảnh báo có thể tra cứu. Giá trị của đề tài nằm ở tổ chức hệ thống quanh các mô hình có sẵn; không phải phát triển một mạng nơ-ron mới.

Kiểm thử chức năng trên dữ liệu tổng hợp cung cấp bằng chứng cho những đường xử lý được mô tả ở chương 5. Các kết quả này không thay thế kiểm thử camera trên thiết bị thật hoặc đo độ chính xác nhận diện. Web di động dùng chung frontend; ứng dụng Flutter còn cần hoàn thiện trước khi được xem là một kênh thi tương đương.

\section{Những vấn đề còn tồn tại}
Danh tính sinh viên vẫn do người dùng khai báo, chưa đối chiếu với hồ sơ đào tạo hoặc nhận diện sinh trắc học. Nhãn cảnh báo được tạo tại client, chưa có cơ chế bảo đảm tín hiệu không bị sửa đổi. Hệ thống chưa có hàng đợi sự kiện bền vững khi mất mạng hoặc lưu nháp câu trả lời trên server.

Ngưỡng AI chưa được hiệu chuẩn bằng bộ dữ liệu độc lập. Một nhãn ưu tiên có thể che dấu hiệu khác trong cùng khung hình. Vòng đời camera/timer sau nộp, độ ổn định trên trình duyệt di động và khả năng chịu tải vẫn cần thử nghiệm bổ sung. Phân quyền theo lớp, quy trình xét duyệt cảnh báo và nhật ký truy cập chưa hoàn chỉnh.

\section{Hướng phát triển}
\begin{longtable}{@{}L{1.3cm}L{6.2cm}L{7.1cm}@{}}
\caption{Các hướng phát triển tiếp theo}\\
\toprule Ưu tiên & Nội dung & Tiêu chí đánh giá \\\midrule\endhead
1 & Xác minh danh tính theo danh sách lớp; kiểm tra phân quyền nhất quán; hoàn thiện kết thúc phiên và camera. & Phiên đúng đối tượng; dữ liệu đúng phạm vi; tài nguyên dừng sau nộp.\\
2 & Lưu nháp server; hàng đợi sự kiện có retry và ID chống trùng; cải thiện phục hồi khi phản hồi cấp phiên bị mất. & Khôi phục theo chính sách; không mất hoặc nhân đôi dữ liệu ngoài dự kiến.\\
3 & Thu dữ liệu có nhãn, hiệu chỉnh ngưỡng và đo hiệu năng thiết bị. & Precision, recall, F1, cảnh báo nhầm/giờ và độ trễ theo từng nhãn.\\
4 & Đồng bộ Flutter với phiên thi; thay phân loại mô phỏng bằng mô hình phù hợp. & Luồng đầu cuối trên thiết bị thật và bằng chứng nhận diện.\\
5 & Hoàn thiện migration có phiên bản, kiểm thử tải và quy trình khôi phục. & Tái lập môi trường, nâng cấp dữ liệu và đo giới hạn trên cấu hình xác định.\\
\bottomrule
\end{longtable}

